
\begin{table}[t]
\centering
\small
\caption{Locked validation protocol used for the primary deployable analysis. This table separates the unbiased locked analysis from the diagnostic upper-bound envelope.}
\label{tab:locked-protocol}
\begin{threeparttable}
\begin{tabularx}{\textwidth}{p{0.24\textwidth}X}
\toprule
Component & Protocol \\
\midrule
Split design & For each model--dataset--seed condition, examples are split stratified by correctness into 60\% train, 20\% validation, and 20\% held-out test. Five seeds are used. \\
Feature fitting & Scalers and PCA transformations are fitted on the train split only. Validation and test examples are transformed using train-fitted objects. \\
Probe families & Logistic reliability selectors are trained for confidence/option features, PCA-hidden features, and confidence/option plus PCA-hidden features. The logistic regularization grid is $C\in\{0.1,1,10\}$. \\
Validation-locked selection & Probe variant, PCA dimension ($16,32,64$), and regularization strength are selected using validation AUROC within each feature family. \\
Protected test reporting & AUROC, AUPRC-failure, Risk@coverage, AURC, E-AURC, and failure capture are reported once on the held-out test split after validation selection. \\
Diagnostic envelope & The broader best-probe envelope in later tables is retained only as an exploratory upper-bound diagnostic; it is not used as the primary deployable estimate. \\
Uncertainty & Primary tables report bootstrap intervals over model--dataset conditions; paired deltas use paired condition--seed comparisons with bootstrap intervals and Wilcoxon signed-rank tests. \\
\bottomrule
\end{tabularx}
\begin{tablenotes}[flushleft]\footnotesize\item This locked protocol addresses optimistic selection risk: richer hidden-augmented feature families are not allowed to choose probes or PCA dimensions using the final test fold.\end{tablenotes}
\end{threeparttable}
\end{table}
